% §3 Ordinary-Bench: Data & Protocol
% Target: ~1.5 pages for NeurIPS Datasets & Benchmarks track

\section{Ordinary-Bench: Data \& Protocol}
\label{sec:benchmark}

Ordinary-Bench is a controlled benchmark for evaluating VLM spatial understanding through dense ordinal probing. It comprises 700 procedurally generated 3D scenes rendered in Blender, two families of ordinal spatial tasks (QRR and TRR) yielding over 336{,}000 probes in total, and a three-condition evaluation protocol that separates visual perception from language priors. This section describes each component.

\subsection{Scene Generation}
\label{sec:scene-gen}

Scenes are generated using a Blender-based rendering pipeline. Each scene places $n$ geometric primitives (spheres, cubes, cylinders) on a ground plane with randomized positions, sizes (\texttt{small}, \texttt{medium}, \texttt{large}), materials (\texttt{metal}, \texttt{rubber}), colors, and rotations. A minimum inter-object distance constraint ($d_{\min} = 0.25$ Blender units) prevents overlap. Lighting and camera parameters are fixed across scenes to control for visual confounds: the primary camera is positioned at azimuth $45^\circ$, elevation $30^\circ$, distance $12.0$ units, rendering $480 \times 320$ images at 256 Cycles samples. Each scene additionally provides four viewpoints for multi-view experiments.

We generate seven \emph{complexity splits}, indexed by object count $n \in \{4, 5, \ldots, 10\}$. Each split contains 100 scenes produced from distinct random seeds, yielding $7 \times 100 = 700$ scenes in total. This design provides a controlled complexity gradient: the number of spatial relations grows combinatorially with $n$, letting us study how model performance degrades as scenes become more crowded. For each scene, the pipeline outputs (i)~a rendered PNG image, (ii)~a JSON file containing each object's 3D coordinates, pixel coordinates, depth, and categorical attributes (shape, color, size, material), and (iii)~multi-view images from four camera positions.

\subsection{Task Definition}
\label{sec:tasks}

We define two families of ordinal spatial tasks that together probe both metric (distance) and directional (bearing) aspects of spatial understanding. Both are defined over ground-truth 3D coordinates, making them view-invariant.

\paragraph{Quaternary Relative Relations (QRR).}
A QRR probe presents four objects organized as two pairs, $(A, B)$ and $(C, D)$, and asks: \emph{Is the 3D distance between $A$ and $B$ less than, approximately equal to, or greater than the distance between $C$ and $D$?} The model must respond with one of three comparators: $<$, $\approx$, or $>$. Approximate equality is determined by a tolerance parameter $\tau = 0.10$: we define $d(A,B) \approx_\tau d(C,D)$ iff $|d(A,B) - d(C,D)| \leq \tau \cdot \max(d(A,B), d(C,D))$, with strict inequality otherwise. We restrict to \emph{disjoint} pairs (no shared objects) to avoid trivially correlated judgments, and filter out boundary cases where the raw metric difference lies between $0.8\tau$ and $1.2\tau$ times the maximum distance (i.e., near the decision boundary) to ensure reliable labels. For $n$ objects, the number of disjoint QRR probes is $\binom{n}{2}\binom{n-2}{2}/2$ before boundary filtering; for instance, a 10-object scene yields up to 630 QRR probes.

\paragraph{Ternary Relative Relations (TRR).}
A TRR probe specifies three objects---a \emph{target}, a \emph{reference origin} (\texttt{ref1}), and a \emph{reference direction} (\texttt{ref2})---and asks: \emph{Standing at \texttt{ref1} and facing toward \texttt{ref2} (the 12~o'clock direction), at what clock hour (1--12) does the target appear?} The ground-truth answer is computed by projecting the target's position onto the reference coordinate frame:
\begin{equation}
  \alpha = \bigl(\operatorname{atan2}(y_t - y_1,\; x_t - x_1) - \operatorname{atan2}(y_2 - y_1,\; x_2 - x_1)\bigr) \bmod 360^\circ,
  \label{eq:trr-angle}
\end{equation}
where subscripts $t$, $1$, $2$ denote target, \texttt{ref1}, \texttt{ref2} coordinates respectively. The angle $\alpha$ is then discretized into clock hours, each spanning a $30^\circ$ sector. Evaluation uses two granularities: \emph{hour-exact} (correct hour) and \emph{quadrant-correct} (correct $90^\circ$ quadrant). For $n$ objects, TRR probes enumerate all ordered triples: $P(n, 3) = n(n{-}1)(n{-}2)$; a 10-object scene yields 720 TRR probes.

\paragraph{Probe counts.}
Table~\ref{tab:probe-counts} summarizes the probe counts across complexity splits. The total benchmark contains 138{,}600 QRR probes and 197{,}400 TRR probes (before boundary filtering), yielding approximately 336{,}000 ordinal probes across all 700 scenes. The combinatorial growth from $n{=}4$ (27 probes/scene) to $n{=}10$ (1{,}350 probes/scene) provides a natural stress test of model capacity.

\begin{table}[t]
\centering
\small
\caption{Probe counts per scene and per split (100 scenes each). QRR counts are before boundary filtering; actual counts are slightly lower due to tolerance-boundary exclusion.}
\label{tab:probe-counts}
\begin{tabular}{c r r r r}
\toprule
$n$ & QRR/scene & TRR/scene & Total/scene & Split total \\
\midrule
4  &     3 &    24 &      27 &     2{,}700 \\
5  &    15 &    60 &      75 &     7{,}500 \\
6  &    45 &   120 &     165 &    16{,}500 \\
7  &   105 &   210 &     315 &    31{,}500 \\
8  &   210 &   336 &     546 &    54{,}600 \\
9  &   378 &   504 &     882 &    88{,}200 \\
10 &   630 &   720 &   1{,}350 &   135{,}000 \\
\midrule
\textbf{Total} & \multicolumn{3}{r}{} & \textbf{336{,}000} \\
\bottomrule
\end{tabular}
\end{table}

\subsection{Evaluation Protocol}
\label{sec:protocol}

\paragraph{Batch JSON prompting.}
For each scene, probes are grouped into batches (default size 20) and presented to the VLM alongside the rendered image and a structured list of object descriptions. The model is instructed to return a JSON array of \texttt{\{qid, answer\}} pairs. A ReAct-style correction loop re-prompts the model if the initial response contains parsing errors or missing answers. This batch format maximizes throughput while ensuring structured, machine-parseable outputs.

\paragraph{Three-condition protocol.}
To disentangle visual perception from language-driven guessing, we evaluate each model under three conditions on identical question sets:
\begin{itemize}[nosep]
  \item \textbf{Condition~A} (correct image): the scene's own rendered image is provided---the standard VQA setup.
  \item \textbf{Condition~B} (mismatched image): a randomly selected image from a \emph{different} scene is provided, so the visual input is valid but irrelevant to the questions.
  \item \textbf{Condition~C} (no image): no image is provided; the model receives only textual object descriptions.
\end{itemize}
Comparing performance across conditions reveals the degree to which a model's spatial answers are grounded in the visual input. We define the \emph{Visual Information Gain} (VIG) as the improvement in reconstruction quality (measured by NRMS; see \S\ref{sec:recon-metrics}) from Condition~C to Condition~A: $\mathrm{VIG} = \mathrm{NRMS}_C - \mathrm{NRMS}_A$. A positive VIG indicates that the correct image genuinely improves the model's spatial belief; a near-zero VIG suggests the model relies primarily on language priors regardless of visual input.

\paragraph{Scoring.}
QRR probes are scored by exact comparator match. TRR probes are evaluated at two granularities: \emph{hour-exact} (predicted hour matches ground truth) and \emph{quadrant-correct} (predicted hour falls in the same $90^\circ$ quadrant as the ground truth). We additionally track \emph{adjacent} accuracy (predicted hour within $\pm 1$ of ground truth, wrapping at 12$\leftrightarrow$1) as a lenient measure. All metrics are reported per-split and aggregated.

\subsection{Dataset Statistics and Release}
\label{sec:dataset-stats}

Ordinary-Bench contains 700 scenes across 7 complexity splits, with 100 scenes each. Each scene includes a single-view rendered image, four multi-view images, and a JSON annotation file specifying object attributes and 3D/2D coordinates. The total dataset comprises approximately 336{,}000 ordinal probes (138{,}600 QRR $+$ 197{,}400 TRR before boundary filtering). Scenes are partitioned into train and test subsets within each split to support model development without test-set contamination.

Object attributes span 3 shapes (sphere, cube, cylinder), 3 sizes (small, medium, large), 2 materials (metal, rubber), and 8 colors, placed on a ground plane ($z = 0$ for all objects). The procedural generation ensures full coverage of the complexity gradient while maintaining exact geometric ground truth. The release package is intended to include the dataset artifacts, generation code, evaluation scripts, question generation pipeline, and accompanying dataset documentation so that the full benchmark can be reproduced and extended.
